\section{Introducción}
\label{sec:introduccion}

El objetivo de este proyecto es el desarrollo de un sistema inteligente capaz de extraer automáticamente el importe total de una factura o recibo a partir de su imagen. En un contexto donde la digitalización de documentos financieros avanza a gran velocidad, la capacidad de interpretar automáticamente estos documentos se convierte en una necesidad técnica para automatizar procesos contables, de auditoría y de gestión empresarial. El sistema propuesto actúa como un puente entre la imagen sin estructura que captura un dispositivo y el dato numérico preciso que necesita un sistema de gestión.

El núcleo de la investigación se centra en la combinación de dos tecnologías complementarias: el reconocimiento óptico de caracteres (OCR), que convierte la imagen en texto, y un modelo de aprendizaje automático que, a partir del texto extraído, identifica y selecciona el candidato numérico que corresponde al total de la factura. A diferencia de los enfoques basados exclusivamente en expresiones regulares o búsqueda de palabras clave, que fallan ante la variabilidad de formatos y los errores inherentes al OCR, el sistema emplea un clasificador Gradient Boosting entrenado sobre un conjunto de características diseñadas para capturar el contexto posicional y semántico de cada candidato numérico.

Finalmente, el proyecto destaca por su orientación a la eficiencia y el despliegue en dispositivos con recursos limitados. A lo largo de la experimentación se han evaluado ocho modelos de extracción y dos motores OCR, priorizando no solo la precisión, sino también el tiempo de inferencia y el tamaño del modelo resultante. El pipeline final, basado en PaddleOCR y Gradient Boosting, alcanza un 83,2\% de coincidencia exacta con un tiempo de procesamiento inferior a un segundo por imagen, lo que lo convierte en una solución viable para aplicaciones móviles y en tiempo real.
